\documentclass[11pt,a4paper]{article}
\usepackage[margin=1in]{geometry}
\usepackage{graphicx}
\usepackage{float}
\usepackage{hyperref}
\usepackage{booktabs}
\usepackage{caption}
\usepackage{subcaption}
\usepackage{enumitem}

\setlength{\parskip}{0.35em}
\setlength{\parindent}{0pt}
\setlist[itemize]{nosep,leftmargin=1.5em}
\setlength{\textfloatsep}{8pt plus 1pt minus 2pt}
\setlength{\floatsep}{8pt plus 1pt minus 2pt}
\setlength{\intextsep}{8pt plus 1pt minus 2pt}
\captionsetup{skip=4pt,font=small}

\title{Part 1 Report: Traditional Filters Implementation}
\author{}
\date{\today}

\begin{document}
\maketitle

\section{Introduction}
This report summarizes the testing and evaluation for Part 1 of the assignment: implementing classical image filtering methods from scratch.

Implemented components:
\begin{itemize}
    \item 2D convolution
    \item Mean filter
    \item Gaussian filter and Gaussian kernel generation
    \item Laplacian filter (standard and diagonal variants)
    \item Sobel filter (x, y, and magnitude/direction)
    \item Bonus: Canny edge detector
\end{itemize}

\section{Methodology}
The implementation was tested on four required image categories:
\begin{itemize}
    \item Synthetic image with geometric structures: \texttt{example\_images/Sythetic.png}
    \item Natural image with textures: \texttt{example\_images/natural.jpg}
    \item Portrait/face image: \texttt{example\_images/Face.png}
    \item Noisy image (Gaussian + salt-and-pepper): \texttt{example\_images/noise.jpg}
\end{itemize}

For each image, the following were generated:
\begin{itemize}
    \item Basic filter visualization
    \item Comparison with OpenCV for smoothing and edge operators
    \item Canny comparison grid (bonus)
\end{itemize}

Outputs were saved in \texttt{filter\_results/}.

\section{Testing Results}

\subsection{Synthetic Image (Geometric)}
\begin{figure}[H]
    \centering
    \begin{subfigure}[t]{0.49\linewidth}
        \centering
        \includegraphics[width=\linewidth]{filter_results/Sythetic_basic_filters.png}
        \caption{Basic filtering results.}
    \end{subfigure}
    \hfill
    \begin{subfigure}[t]{0.49\linewidth}
        \centering
        \includegraphics[width=\linewidth]{filter_results/Sythetic_canny_comparison.png}
        \caption{Canny comparison (ours vs OpenCV).}
    \end{subfigure}
    \vspace{4pt}
    \begin{subfigure}[t]{0.49\linewidth}
        \centering
        \includegraphics[width=\linewidth]{filter_results/Sythetic_opencv_comparison_smoothing.png}
        \caption{Smoothing comparison with OpenCV.}
    \end{subfigure}
    \hfill
    \begin{subfigure}[t]{0.49\linewidth}
        \centering
        \includegraphics[width=\linewidth]{filter_results/Sythetic_opencv_comparison_edges.png}
        \caption{Edge operator comparison with OpenCV.}
    \end{subfigure}
    \caption{Synthetic image results.}
\end{figure}

Observations:
\begin{itemize}
    \item Edges are sharp and highly structured, so Sobel and Laplacian clearly highlight boundaries.
    \item Mean/Gaussian smoothing noticeably softens checker-like transitions.
    \item Our outputs and OpenCV outputs are visually very close.
\end{itemize}

\subsection{Natural Image (Texture-Rich)}
\begin{figure}[H]
    \centering
    \begin{subfigure}[t]{0.49\linewidth}
        \centering
        \includegraphics[width=\linewidth]{filter_results/natural_basic_filters.png}
        \caption{Basic filtering results.}
    \end{subfigure}
    \hfill
    \begin{subfigure}[t]{0.49\linewidth}
        \centering
        \includegraphics[width=\linewidth]{filter_results/natural_canny_comparison.png}
        \caption{Canny comparison (ours vs OpenCV).}
    \end{subfigure}
    \vspace{4pt}
    \begin{subfigure}[t]{0.49\linewidth}
        \centering
        \includegraphics[width=\linewidth]{filter_results/natural_opencv_comparison_smoothing.png}
        \caption{Smoothing comparison with OpenCV.}
    \end{subfigure}
    \hfill
    \begin{subfigure}[t]{0.49\linewidth}
        \centering
        \includegraphics[width=\linewidth]{filter_results/natural_opencv_comparison_edges.png}
        \caption{Edge operator comparison with OpenCV.}
    \end{subfigure}
    \caption{Natural image results.}
\end{figure}

Observations:
\begin{itemize}
    \item Gaussian filtering preserves large structures while suppressing high-frequency texture.
    \item Sobel magnitude captures road, object contours, and dominant texture gradients.
    \item Canny with larger sigma produces cleaner but less detailed edge maps.
\end{itemize}

\subsection{Portrait/Face Image}
\begin{figure}[H]
    \centering
    \begin{subfigure}[t]{0.49\linewidth}
        \centering
        \includegraphics[width=\linewidth]{filter_results/Face_basic_filters.png}
        \caption{Basic filtering results.}
    \end{subfigure}
    \hfill
    \begin{subfigure}[t]{0.49\linewidth}
        \centering
        \includegraphics[width=\linewidth]{filter_results/Face_canny_comparison.png}
        \caption{Canny comparison (ours vs OpenCV).}
    \end{subfigure}
    \vspace{4pt}
    \begin{subfigure}[t]{0.49\linewidth}
        \centering
        \includegraphics[width=\linewidth]{filter_results/Face_opencv_comparison_smoothing.png}
        \caption{Smoothing comparison with OpenCV.}
    \end{subfigure}
    \hfill
    \begin{subfigure}[t]{0.49\linewidth}
        \centering
        \includegraphics[width=\linewidth]{filter_results/Face_opencv_comparison_edges.png}
        \caption{Edge operator comparison with OpenCV.}
    \end{subfigure}
    \caption{Face image results.}
\end{figure}

Observations:
\begin{itemize}
    \item Facial boundary and clothing folds are preserved by Sobel/Canny.
    \item Gaussian filter reduces skin-level noise more naturally than mean filter.
    \item Laplacian diagonal kernel shows slightly stronger response around fine details.
\end{itemize}

\subsection{Noisy Image}
\begin{figure}[H]
    \centering
    \begin{subfigure}[t]{0.49\linewidth}
        \centering
        \includegraphics[width=\linewidth]{filter_results/noise_basic_filters.png}
        \caption{Basic filtering results.}
    \end{subfigure}
    \hfill
    \begin{subfigure}[t]{0.49\linewidth}
        \centering
        \includegraphics[width=\linewidth]{filter_results/noise_canny_comparison.png}
        \caption{Canny comparison (ours vs OpenCV).}
    \end{subfigure}
    \vspace{4pt}
    \begin{subfigure}[t]{0.49\linewidth}
        \centering
        \includegraphics[width=\linewidth]{filter_results/noise_opencv_comparison_smoothing.png}
        \caption{Smoothing comparison with OpenCV.}
    \end{subfigure}
    \hfill
    \begin{subfigure}[t]{0.49\linewidth}
        \centering
        \includegraphics[width=\linewidth]{filter_results/noise_opencv_comparison_edges.png}
        \caption{Edge operator comparison with OpenCV.}
    \end{subfigure}
    \caption{Noisy image results.}
\end{figure}

Observations:
\begin{itemize}
    \item Both mean and Gaussian reduce noise.
    \item Gaussian denoising preserves edges slightly better than mean filtering.
    \item In very noisy regions, Canny thresholding requires careful tuning to avoid fragmented edges.
\end{itemize}

\section{Parameter Analysis}
\subsection{Kernel Size and Sigma}
\begin{itemize}
    \item Larger kernel size in mean/Gaussian gives stronger smoothing and more detail loss.
    \item Larger sigma in Gaussian and Canny pre-smoothing suppresses texture/noise but can remove weak edges.
\end{itemize}

\subsection{Sobel Directionality}
\begin{itemize}
    \item Sobel-X emphasizes vertical boundaries.
    \item Sobel-Y emphasizes horizontal boundaries.
    \item Magnitude combines both and is more suitable for generic edge intensity visualization.
\end{itemize}

\subsection{Canny Thresholds}
From the Canny comparison grids:
\begin{itemize}
    \item Lower thresholds produce denser edges but include more noise.
    \item Higher thresholds produce cleaner maps but may miss weak contours.
\end{itemize}

\section{Border Handling Analysis}
\begin{figure}[H]
    \centering
    \includegraphics[width=0.82\linewidth]{filter_results/padding_mean_comparison.png}
    \caption{Mean filter border handling comparison: constant, reflect, and replicate padding.}
\end{figure}

Interpretation:
\begin{itemize}
    \item Constant padding tends to darken borders and may introduce artificial edge responses near boundaries.
    \item Reflect padding preserves continuity and usually gives the most natural border behavior.
    \item Replicate padding keeps border intensity constant and often performs similarly to reflect for smooth regions.
\end{itemize}

\section{Noisy vs Clean Analysis}
\begin{itemize}
    \item Both mean and Gaussian decrease intensity variation (noise suppression effect).
    \item Gaussian retains structure slightly better visually while still reducing random noise.
\end{itemize}

\section{Conclusion}
Part 1 implementation is correct and complete with required filters and bonus Canny.

Overall conclusions:
\begin{itemize}
    \item Functional correctness is validated by visual agreement with OpenCV across multiple image categories.
    \item Parameter settings strongly control smoothing-strength and edge sensitivity.
    \item Border strategy affects edge artifacts near image boundaries; reflect/replicate are typically preferable to constant padding.
    \item For noisy inputs, Gaussian filtering provides a better blur-detail trade-off than mean filtering.
\end{itemize}

\end{document}
